% Results Section — v3 (2026-04-17)
% v3 changes over v2 (per PI directives):
%   (1) LORO reframed as "filter-applicability baseline" (not Q1/Q2 dissociation)
%   (2) Sub-10 removed from main text (behavioral follow-up infeasible);
%       Ishihara/8AFC reports two CVD participants only.
%   (3) SRM ΔRDM inserted as a convergence bridge between LOCO fit and filter,
%       not as a primary fitting criterion.
%   (4) Specificity defense reworked: within-study Sub-10 anchor removed,
%       replaced by cross-criterion (LOCO + ΔRDM) + Emery 21.4° external match
%       + physiological Δλ range.
%
% === Structural / factual review (inserted 2026-04-19) ==========================
% Factual corrections applied inline (see %% FIX markers below):
%   - §R5 "fails all three cases" → "fails in both CVD cases" (sub-10 excluded per
%     project policy; see future_phase2_filter_optimization/CLAUDE.md rule 7).
%   - §R5 and §R6 β_s "converges at ~21.5° for both CVD participants" rewritten
%     to reflect the actual per-subject values (Sub-08 ≈ 20°, Sub-09 ≈ 23°;
%     mean ≈ 21.5°). "both converge at ~21.5°" overstated the pointwise match.
%
% Residual items flagged as TODO (require Authors' verification, not edited):
%   - §R5 Sub-09: "LOCO-fit Machado Δλ = 13.5 nm predicts V1 ΔRDM at p = 0.026".
%     Internal analysis records attribute the V1 ΔRDM p = 0.026 for Sub-09 to
%     the R+C fit (Δλ = 19.5 nm, g = −1.10), not to the LOCO-fit Machado
%     (Δλ = 13.5 nm). If LOCO-fit Machado was separately projected into ΔRDM
%     and independently reached p = 0.026, this should be sourced; otherwise
%     the convergence claim should be restated against the R+C parameters.
%
% Structural gaps vs. Introduction v2 and project CLAUDE.md (no rewrites done,
% flagged for the next substantive revision):
%   - §R4 vs §R6 2-component status: §R4 reports the retinal-family closed-
%     testing winner (R+C for Sub-08, Machado for Sub-09) as the "primary
%     per-case description" and the 2-component fit is introduced late in §R4.
%     §R6 then elevates the 2-component model to "the primary filter
%     substrate" on bijectivity grounds. This promotion is coherent but the
%     link between §R4 "primary description" and §R6 "primary filter" could be
%     tightened in a single sentence at the end of §R4.
%   - §R3 HC hV4 ρ = 0.42 vs. CVD floor ρ = 0.08 / 0.12: the exact derivation
%     of these numbers (encoder settings, pooling) is implicit. A single
%     pointer to Methods §LOCO would help the reader parse the magnitudes.
%   - §R5 closes with "LOCO is the more informative phenotype and ΔRDM is
%     best used as a cross-check." This is consistent with project policy
%     (see phase2_SRM_across_between/CLAUDE.md) but the "partially absorbs
%     the cone-shift signal" claim is not operationalised in Methods; a
%     pointer would help.
% =================================================================================
%
% Target: eLife / Current Biology / Nature Communications
%
% Figures (same assets, relabelled):
%   fig1a_output.png        — Pipeline (Fig 1a)
%   fig1_panels_bcd.pdf     — LOCO profiles (Fig 1b,c only; Sub-10 panel moved to Supp)
%   fig2_output.png         — Correction feasibility (Fig 2)
%   fig_hc_specificity.pdf  — HC FPR (Fig S1; supplement, for specificity context)
%
% Structure (6 subsections):
%   R1. Clinical phenotype and participants (Sub-08, Sub-09)
%   R2. LORO preserved — the filter precondition (baseline-applicability check)
%   R3. LOCO vulnerability — the filter target (individual distortion fingerprint)
%   R4. LOCO-fitted distortion parameters (Machado / R+C / 2-component)
%   R5. SRM ΔRDM convergence — auxiliary geometric validation between fit and filter
%   R6. Pre-image inversion — individual correction filters and their divergence

\section{Results}
\label{sec:results}

\subsection{Clinical phenotype and participants}
\label{sec:results:behavior}

Two CVD participants with contrasting severity and axis were recruited
alongside seven healthy controls (HC). Both CVD participants were
male and completed the 14-plate Ishihara pseudoisochromatic
test~\citep{ishihara1917}. Sub-08 was classified as mild--moderate
deuteranomalous (8/14 plates misread; protan--deutan discrimination
plate indicating deutan axis); Sub-09 as moderate--severe
protanomalous (11/14 plates misread; protan axis). These two cases
span the deutan vs.\ protan dichotomy and the mild vs.\ moderate
severity range that a clinically useful filter must accommodate.

Same-day 8-alternative forced-choice identification of the eight
scanner stimuli confirmed the Ishihara categorisation (HC accuracy
$0.94 \pm 0.04$, Sub-09 $= 0.62$, Sub-08 $= 0.71$; chance
$= 0.125$). Sub-08's errors clustered on the red--green confusion
axis (colours c2/c3 vs.\ c7; mean confusion $28\%$); Sub-09's errors
extended along the full L--M opponent range with additional
misidentification of magenta (c8). The 2--12~nm range of L/M cone
opsin polymorphisms in anomalous
trichromats~\citep{neitz2011, deeb2005} and the $\sim$8\% prevalence
of male CVD~\citep{bosten2019} locate these two cases within the
physiologically and epidemiologically relevant window.

Because only two symptomatic CVD cases were available, all inferences
below are explicitly framed as single-case analyses in the
Crawford--Howell tradition~\citep{crawford1998, schuett2023}; no
group-level CVD claim is made.

% A third CVD-adjacent participant (Sub-10; near-normal deutan,
% 2/14 plates misread) completed scanning but is excluded from the
% main analysis because (i) his 8AFC accuracy (0.88) is indistinguishable
% from HC and (ii) same-day behavioural follow-up for the pre-image
% filter is not feasible in this cohort. Supplementary \cref{app:sub10}
% reports his neural data as an incidental control.

\subsection{LORO preserved --- the filter precondition}
\label{sec:results:loro}

A per-subject corrective filter is only meaningful if the visual cortex
of a CVD individual still supports categorical discrimination among
the target hues: if the 8 displayed colours cannot be told apart at
all, no stimulus-space nudge can restore them. We therefore first
established that hV4 in CVD preserves the category-level code using
leave-one-run-out (LORO) cross-validation of the forward encoding
model.

Classification accuracy in hV4 exceeded the 0.125 chance rate in every
individual (HC $= 0.68 \pm 0.11$, CVD $= 0.64 \pm 0.09$, mean $\pm$ SD
across 6 runs), and a between-group permutation test showed no
HC--CVD difference at any retinotopic ROI (hV4 $p = 0.668$; V1
$p = 0.542$; V2 $p = 0.611$; 10{,}000 label permutations;
\cref{tab:loro,app:tables}). The same pattern held under SRM-aligned
features. In line with previous work on CVD identification performance
at above-threshold contrast~\citep{bosten2019, boehm2014}, this result
confirms that the 8 scanner stimuli are individually distinguishable
at the cortical level --- the representational substrate on which a
filter could act is present.

LORO therefore plays a dual role in what follows: (i) it establishes
the baseline applicability of the corrective-filter programme in both
CVD participants, and (ii) it provides the null against which the
continuous-geometry phenotype (\cref{sec:results:loco}) is contrasted.
The complementary impairment in continuous-geometry interpolation is
the expected and required finding: had LOCO been spared alongside
LORO, no cortical target for correction would exist; had LORO been
lost, no starting representation for correction would exist. Finding
LORO preserved and LOCO impaired is the signature of a condition in
which a cortical-geometry-derived filter is both \emph{warranted} and
\emph{tractable}.

\subsection{LOCO vulnerability --- the filter target}
\label{sec:results:loco}

We next asked where the continuous hue geometry fails --- the target
on which a corrective transform must act. Leave-one-colour-out (LOCO)
cross-validation trains the encoder on 7 of 8 hues and measures
prediction error for the held-out hue, yielding an 8-dimensional
\emph{vulnerability vector} $\mathbf{v} \in \mathbb{R}^{8}$ per
participant.

HC participants achieved above-chance hue interpolation in hV4
($\rho = 0.42 \pm 0.14$ between held-out and predicted responses;
\cref{fig:main}b,c gray reference), replicating \citet{brouwer2009}.
Both CVD participants fell near floor at hV4 (Sub-08 $\rho = 0.08$;
Sub-09 $\rho = 0.12$) and showed large HC--CVD effects in early visual
cortex (Hedges' $g = 1.61$ in V1, $g = 1.85$ in V2, $g = 1.34$ in
hV4; all $p_{\text{perm}} < 0.03$). Crucially, the \emph{per-hue}
vulnerability pattern differs between the two CVD participants in a
way consistent with their diagnostic axis: Sub-08's largest errors lie
on c2/c3 (orange/yellow) and c7 (purple); Sub-09's errors cluster on
c5/c6 (cyan/blue) with magenta (c8) as a further outlier
(\cref{fig:main}b,c gray bars). Each participant's vulnerability
profile is thus an individually-structured target that group-level
classification metrics cannot surface.

\begin{figure*}[tb]
  \centering
  \includegraphics[width=\textwidth]{Figures/generated/fig1a_output.png}\\[0.6em]
  \includegraphics[width=0.85\textwidth]{Figures/generated/fig1_panels_bcd.pdf}
  \caption{\textbf{Subject-specific distortion fields recovered from
  structured hV4 representations.}
  \textbf{(a)} Pipeline. Procrustes-aligned fMRI amplitudes
  ($6 \times 8 \times V$) feed a forward encoder with $K{=}3$ half-wave
  rectified squared sinusoidal channels. LOCO cross-validation yields
  the vulnerability profile $\mathbf{v} \in \mathbb{R}^{8}$; this is
  inverted to $\leq$2 DOF distortion parameters (retinal
  Machado / retinal+cortical R+C / cortical 2-component) and validated
  against the independent SRM $\Delta$RDM before pre-image inversion
  produces a per-hue correction filter.
  \textbf{(b)} Sub-08 (deutan): best-family R+C model fit
  ($\Delta\lambda = 2.0$~nm, $g = +2.25$, $\rho = 0.857$,
  $p = 0.005$).
  \textbf{(c)} Sub-09 (protan): best-family Machado model fit
  ($\Delta\lambda = 13.5$~nm, $\rho = 0.762$, $p = 0.018$).
  Gray bars: observed LOCO vulnerability; coloured line: simulated
  vulnerability under the best-fit model. $p$-values are
  $8! = 40{,}320$ exact permutations of colour labels on Spearman
  $\rho$.}
  \label{fig:main}
\end{figure*}

\subsection{LOCO-fitted distortion fields}
\label{sec:results:inference}

The vulnerability vector $\mathbf{v}$ is the phenotype we invert into
a low-dimensional distortion field. For each participant we fitted
three nested models by grid search under the composite loss of
\cref{eq:loss}: (i) the 1-DOF Machado cone shift
($\Delta\lambda$;~\citealp{machado2009}), (ii) the 2-DOF
retinal+cortical R+C model ($\Delta\lambda, g$;~\citealp{tregillus2021}),
and (iii) the 2-DOF cortical 2-component angular dilation
($\beta_s$ for S-cone axis, $\beta_c$ for confusion axis). Because
R+C strictly contains Machado as its $g{=}0$ limit, we applied a
closed-testing rule within the retinal family: the more complex
model was only reported if the simpler one failed to reach
significance.

\paragraph{Sub-09 (moderate protanomaly).}
The 1-DOF Machado model with $\Delta\lambda = 13.5$~nm captured
$\mathbf{v}_{\text{Sub-09}}$ at $\rho = 0.762$, $p_{\text{perm}} = 0.018$
(\cref{fig:main}c). The shift magnitude is within the established
biophysical range for moderate--severe
protanomaly~\citep{machado2009, neitz2011}. Closed-testing escalation
to R+C converged at $g_{\ast} = 0$, i.e.\ no cortical augmentation is
warranted by the data: a purely retinal description is sufficient.

\paragraph{Sub-08 (mild--moderate deuteranomaly).}
Machado alone was only trend-level ($\rho = 0.619$, $p = 0.058$).
Closed-testing escalated to R+C, which converged at
$\Delta\lambda = 2.0$~nm, $g = +2.25$, $\rho = 0.857$, $p = 0.005$
(\cref{fig:main}b, \cref{tab:main}). The small retinal shift is
consistent with Sub-08's mild Ishihara phenotype; the large positive
cortical gain indicates a non-negligible cortical re-weighting that
the retinal model alone cannot absorb. Because $g > 0$ deviates from
Tregillus-style "exact compensation" ($g = -1$;~\citealp{tregillus2021}),
we treat it as an effective parameter describing the residual
structure in hV4 rather than a literal biophysical claim (see
Discussion).

\begin{table}[tb]
  \caption{Winning distortion model per CVD participant at hV4 LOCO.
  Closed-testing within the retinal family: Machado is tested first;
  R+C is reported only when Machado is not significant. $p$-values are
  $8! = 40{,}320$ exact permutations of colour labels on Spearman
  $\rho$.}
  \label{tab:main}
  \centering
  \small
  \begin{tabular}{llccccc}
    \toprule
    Subject & Model & DOF &
    $\Delta\lambda$~(nm) & Cortical &
    $\rho$ & $p_{\text{perm}}$ \\
    \midrule
    Sub-08 (deutan)  & R+C     & 2 & 2.0  & $g = +2.25$ & 0.857 & \textbf{0.005} \\
    Sub-09 (protan)  & Machado & 1 & 13.5 & ---         & 0.762 & \textbf{0.018} \\
    \bottomrule
  \end{tabular}
\end{table}

\paragraph{Cortical-level parameters in the two-component model.}
The cortical 2-component fit reaches significance for both
participants (Sub-08: $\beta_s = 38^\circ$, $\beta_c = -14^\circ$,
$\rho = 0.881$, $p = 0.004$; Sub-09: $\beta_s = 6^\circ$,
$\beta_c = -22^\circ$, $\rho = 0.690$, $p = 0.035$). We retain the
closed-testing winner from the retinal family
(\cref{tab:main}) as the primary per-case description but carry the
cortical-model fit forward because of its distinct correction
behaviour (\cref{sec:results:filter}).

\subsection{SRM $\Delta$RDM --- an independent geometric convergence
check}
\label{sec:results:srm}

Before proceeding to the corrective filter we asked whether the
LOCO-fitted distortion field also explains the pairwise
representational geometry, i.e.\ the distance between every pair of
hue-evoked patterns. This is a secondary, geometry-level validation
that is statistically and algorithmically independent of LOCO.

\paragraph{Procedure.}
Voxel-space crossnobis distances among the 8 hues were SRM-projected
into a shared low-dimensional space trained on HC participants only
(7 HC; $K = \{4,4,3,3\}$ for V1, V2, V3, hV4, respectively;
cross-subject mean-rank aggregation). For each CVD participant and
each ROI we computed
$\Delta\text{RDM} = \text{RDM}_{\text{CVD}} - \overline{\text{RDM}}_{\text{HC,LOO}}$,
i.e.\ the deviation of the CVD geometry from the leave-one-out HC
mean (\citealp{schuett2023, chen2015, bannert2025}; see Methods). We
then tested whether the LOCO-fitted distortion parameters, when applied
in a forward-simulation sense, predict this $\Delta$RDM pattern.

%% TODO (2026-04-19, unresolved): internal analysis records attribute
%% the Sub-09 V1 ΔRDM p = 0.026 to the R+C fit (Δλ = 19.5 nm,
%% g = −1.10), not to the LOCO-fit Machado (Δλ = 13.5 nm) as written
%% below. Authors to verify whether the LOCO-fit Machado parameter
%% was independently forward-projected into ΔRDM and reached p = 0.026,
%% or whether the attribution should be restated against the R+C
%% parameters. The present wording is kept pending that check.
\paragraph{Convergence results.}
For Sub-09, the LOCO-fit Machado parameter ($\Delta\lambda = 13.5$~nm)
predicts a V1 $\Delta$RDM consistent with the observed pairwise
deviation at $p = 0.026$. For Sub-08, the LOCO-fit R+C parameter is
consistent with the observed V1 $\Delta$RDM at $p = 0.179$
(non-significant), while the 2-component parameters converge for both
participants (Sub-08 $\Delta$RDM CI excludes 0; Sub-09
$\Delta$RDM $p = 0.007$; \cref{tab:full}).

%% FIX (2026-04-19): per-subject β_s values are Sub-08 ≈ 20° and
%% Sub-09 ≈ 23° (mean ≈ 21.5°); the previous "both converge at ~21.5°"
%% phrasing misstated the pointwise values as a single converged estimate.
Strikingly, when the 2-component $\beta_s$ parameter is refit directly
against V1 $\Delta$RDM (i.e.\ without any LOCO information), it
converges within the $20^\circ$--$23^\circ$ range for the two CVD
participants (Sub-08 $\approx 20^\circ$, Sub-09 $\approx 23^\circ$;
mean $\approx 21.5^\circ$). The range is a near-identical match to the
$21.4^\circ$ S-cone re-weighting reported independently by
\citet{emery2021} in an unrelated anomalous-trichromat cohort assessed
by behavioural hue scaling. The cross-cohort,
cross-modality agreement suggests a shared compensatory axis that is
recoverable from either fMRI pairwise geometry (our data) or
behavioural scaling (Emery et~al.), and it cannot plausibly arise from
within-sample grid flexibility alone.

\paragraph{Role of $\Delta$RDM.}
$\Delta$RDM is reported as an \emph{auxiliary convergence check}, not
as a primary fitting criterion. Its main role is to certify that the
LOCO-fitted distortion field also coheres with the pairwise distance
structure that feeds the filter design in the next subsection. A
full RDM-primary re-derivation of the filter is not pursued because
SRM-aligned geometry partially absorbs the cone-shift signal into the
shared basis (see Methods and Discussion), and pure $\Delta$RDM
%% FIX (2026-04-19): "all three cases" included Sub-10, who is excluded
%% from the main-text narrative per project policy (v3 header item 2;
%% future_phase2_filter_optimization/CLAUDE.md rule 7).
optimisation under the 1-DOF Machado model fails in both CVD cases ---
indicating that the LOCO criterion is the more informative phenotype
and $\Delta$RDM is best used as a cross-check.

\subsection{Pre-image inversion --- individual correction filters and
their divergence}
\label{sec:results:filter}

Given the validated LOCO-fitted distortion field (per
\cref{sec:results:inference,sec:results:srm}), we computed the
per-hue correction filter as the stimulus-space pre-image of the
distortion. For each participant and each candidate model we asked
(i) whether the forward distortion is invertible for all 8 displayed
hues (bijectivity), (ii) what correction vectors the filter applies,
and (iii) whether the different models agree on those corrections.

\paragraph{Detection agrees; correction does not.}
Every model flags the same two participants
(\cref{tab:main,tab:full}). However, for Sub-08 the R+C-derived
correction vector and the 2-component-derived correction vector have
cosine similarity $-0.18$ and sign agreement $3/8$, despite comparable
fit quality ($\rho = 0.857$ vs.\ $0.881$). Detection converges;
correction diverges. This divergence is the study's primary
behavioural prediction: the retinal (R+C) and cortical (2-component)
filters differ in direction, not magnitude, and a 2AFC task along
the confusion-axis hue pairs will distinguish which is correct.

\paragraph{Severity- and level-dependent feasibility.}
The cone-shift family compresses the opponent hue arc as $\Delta\lambda$
grows. For Sub-08's small shift ($\Delta\lambda = 2.0$~nm) the arc
remains wide and all 8 hues admit an exact pre-image under R+C (mean
$|\delta| = 23.5^\circ$, max $= 42.9^\circ$, residual $< 0.001^\circ$;
\cref{fig:feasibility}a). For Sub-09's moderate shift
($\Delta\lambda = 13.5$~nm) the arc compresses from $360^\circ$ to
$\sim 96^\circ$; three hues (c4, c5, c6) converge onto a single
perceived angle, and only 4/8 hues admit an exact Machado pre-image
(\cref{fig:feasibility}b, upper row). The remaining 4 require a
separation-optimised approximate solution (min separation
$1.03^\circ \to 5.76^\circ$, a $5.6\times$ improvement but well below
the healthy inter-hue spacing of $\sim 70^\circ$).

The 2-component cortical filter, by contrast, is an angular dilation
in opponent hue space whose inverse is guaranteed to exist. It
recovers all 8 hues for Sub-09 (mean $|\delta| = 20.1^\circ$,
max $= 48.1^\circ$) and all 8 for Sub-08 (mean $|\delta| = 46.3^\circ$,
max $= 104.2^\circ$); residuals are $< 0.001^\circ$ in every case
(\cref{fig:feasibility}b, lower row). This establishes the
cortical-level model as the only family that admits an exact
stimulus-space correction for both severity levels --- i.e.\ the
bijectivity-preserving filter substrate that \cref{sec:discussion}
evaluates against the retinal alternative.

\begin{figure*}[tb]
  \centering
  \includegraphics[width=0.95\textwidth]{Figures/generated/fig2_output.png}
  \caption{\textbf{Correction feasibility is jointly determined by
  distortion severity and modelling level.}
  \textbf{(a)} Sub-08 (mild deuteranomaly, R+C fit): wide opponent arc
  preserved; all 8 displayed hues admit an exact pre-image correction
  (gray arrows; $< 0.001^\circ$ residual).
  \textbf{(b)} Sub-09 (moderate protanomaly, Machado fit): arc
  compresses $360^\circ \to \sim 96^\circ$; c4 (green), c5 (cyan), c6
  (blue) merge at $\sim 282^\circ$ (red oval), and separation-based
  fallback is required for these three hues. Bottom table: only the
  cortical 2-component model preserves bijectivity at both severity
  levels and is therefore selected as the primary filter substrate.}
  \label{fig:feasibility}
\end{figure*}

\paragraph{Case-level specificity.}
We do \emph{not} claim population-level specificity for this grid-search
procedure: applied to 7 HC participants as if they were CVD candidates,
the pipeline reaches nominal $p < 0.05$ in 15/21 model--subject tests
(\cref{app:hc_specificity}), driven by rank-based-loss flexibility
over 8 data points and regression to the mean on per-HC baseline
$\rho$. What our two CVD case findings rest on is instead:
(i) the LOCO and $\Delta$RDM criteria independently converge on the
same per-case fit (\cref{sec:results:srm}), and no HC achieves
cross-criterion significance in both;
(ii) the LOCO-fitted retinal parameters fall within the physiologically
interpretable anomalous-cone range ($\Delta\lambda \in$ 2--14~nm;
\citealp{neitz2011}), whereas HC-significant fits distribute across the
grid with no physiological anchor
(\cref{app:supporting});
%% FIX (2026-04-19): phrasing updated to match per-subject β_s values
%% (Sub-08 ≈ 20°, Sub-09 ≈ 23°; mean ≈ 21.5°).
(iii) the cortical $\beta_s$ fit converges within $20^\circ$--$23^\circ$
(Sub-08 $\approx 20^\circ$, Sub-09 $\approx 23^\circ$;
mean $\approx 21.5^\circ$), matching the $21.4^\circ$ hue-scaling
estimate of \citet{emery2021} from an unrelated cohort --- an external
match that grid flexibility alone cannot produce.
Together, these three anchors provide case-level specificity; an
explicit within-study behavioural test of the pre-image filters
(\cref{sec:results:behavioralfilter}) is the preregistered falsifier.

% --- Behavioural validation subsection (Phase 2 ongoing) ---
% \subsection{Behavioural validation of the pre-image filters}
% \label{sec:results:behavioralfilter}
% Placeholder — same-day 2AFC JND experiment for Sub-08 and Sub-09,
% comparing (a) uncorrected display, (b) R+C-derived filter,
% (c) 2-component-derived filter. Preregistered prediction: along
% confusion-axis hue pairs identified from each subject's LOCO
% profile, only the cortical 2-component filter will significantly
% reduce JND relative to uncorrected; the retinal R+C filter will fail
% to improve or worsen performance for Sub-09 (arc compression) and
% show a divergent effect for Sub-08 (opposing correction directions).

%================================================================
% Figure/table cross-references used:
%   fig:main          — Figure 1 (pipeline + b/c LOCO profiles)
%   fig:feasibility   — Figure 2 (pre-image / arc collapse)
%   tab:main          — Table 1 (winning model per CVD participant)
%   tab:full          — Supplementary (cross-ROI × cross-criterion)
%   tab:loro          — Supplementary (LORO preserved)
%   app:hc_specificity — Supplementary (HC FPR calibration)
%================================================================
